\subsection{DCL -- Data Control Language}
Die DCL steuert Zugriff und Rechte auf die Datenbank. Kernfrage des Datenschutzes: \emph{Wer darf welche Daten zu welchem Zweck sehen oder ändern?}

\begin{merksatz}[Grundsatz]
Jeder User bekommt nur die Rechte, die er wirklich benötigt -- kein pauschales \icode{ALL}.
\end{merksatz}

\icode{GRANT} vergibt Rechte, \icode{REVOKE} entzieht sie. Rechte sind u.\,a. \icode{ALL}, \icode{SELECT}, \icode{INSERT}, \icode{UPDATE}, \icode{DELETE}, \icode{DROP}. Es gibt drei Ebenen: global (\icode{*.*}), Datenbank (\icode{db.*}) und Tabelle (\icode{db.tabelle}).

\begin{syntaxbox}[SQL]
-- Nutzer anlegen und loeschen
CREATE USER 'test'@'localhost' IDENTIFIED BY 'pw123';
DROP USER   'test'@'localhost';

-- Rechte vergeben und entziehen
GRANT SELECT ON kurs_db.kunde TO 'test'@'localhost';
GRANT ALL    ON kurs_db.*     TO 'test'@'localhost';
REVOKE ALL   ON kurs_db.mitarbeiter FROM 'test'@'localhost';
\end{syntaxbox}
